\section{Bên trong 128 bài, và bên ngoài chúng}
\label{sec:ch12-nghien-cuu-khong-co-nguoi}

\index{nghiên cứu người dùng!chất lượng của phép kiểm chứng|(}%
Mục trước khép lại ở chỗ 128 là chặn trên. Bài báo dành hai mục thảo luận để
nói vì sao, một mục về những bài có người tham gia và một mục về những bài
không có, và hai mục ấy là phần lập luận của bài chứ không phải phần đếm.

Về những bài có người, các tác giả nêu ba chỗ mà một cuộc kiểm chứng có thể
rỗng. Chỗ thứ nhất là ai được hỏi. Trong 128 bài, 13 bài không ghi số người
tham gia và 3 bài chỉ ghi con số ước lượng; nhiều bài không công bố chuyên
môn, trình độ học vấn, hay quan hệ giữa người tham gia với chính nhóm tác giả.
Các tác giả đặt câu hỏi thẳng: nếu không biết gì về nền tảng của những người
này, cũng không biết động cơ nào có thể khiến họ đánh giá một kỹ thuật XAI là
tốt, thì có thể lập luận rằng phép kiểm chứng đã bị lệch một cách không công
bằng~\autocite{p24humangap}.

Chỗ thứ hai là hỏi ai không đúng đối tượng. Nhiều bài tuyển sinh viên đại học
làm người tham gia, trong khi người dùng cuối của kỹ thuật đang được thử là
chuyên gia an ninh mạng, phi công hay bác sĩ. Các tác giả thừa nhận rằng gom
được chuyên gia trong lĩnh vực là việc khó, nhất là những chuyên gia không làm
việc cùng nhóm nghiên cứu, và họ ghi nhận một số bài chỉ có hai hoặc ba chuyên
gia trong cả cuộc đánh giá~\autocite{p24humangap}.

Chỗ thứ ba là hỏi cái gì, và đây là chỗ nặng nhất trong ba. Bài báo chép lại
hai dạng câu hỏi mà họ gặp trong tài liệu. Dạng thứ nhất hỏi thẳng người tham
gia rằng họ có thấy cái này giải thích được hay không, \enquote{Do you think
this is explainable}, và không cho họ chỗ nào để đối chiếu, nên người trả lời
không có cách nào kiểm tra được câu trả lời của mình. Dạng thứ hai là
\enquote{Which is more explainable: our explainable method, or this other
baseline explanation that is clearly worse (e.g., contains grammatical errors
or has logical flaws)?}~\autocite{p24humangap}, tức đặt phương pháp của tác giả
cạnh một mốc so sánh đã hỏng sẵn, nên kết quả biết trước. Nhận xét của bài báo về chỗ này là
một câu về thiết kế thí nghiệm, không phải một lời buộc tội: cần dựng tình
huống đánh giá sát với ngữ cảnh sử dụng thật mà kỹ thuật nhắm
tới~\autocite{p24humangap}. Ba chỗ ấy cộng lại là lý do 128 không nên đọc như
số bài đã kiểm chứng xong.

\index{nghiên cứu người dùng!lập luận tín hiệu}%
Về những bài không có người, các tác giả đưa một lập luận mà chương này lấy
làm trục, và nó đáng chép nguyên. Nếu coi việc tạo ra một sản phẩm có thể giải
thích được như việc phát đi một tín hiệu, thì, các tác giả viết, phải có
xác nhận rằng tín hiệu đã được nhận và được hiểu đúng thì mới xác nhận được
lời tuyên bố về khả năng giải thích; nếu không thì không có gì phân biệt AI
\enquote{giải thích được} với bất kỳ loại AI nào khác, và vô số phương pháp có
thể giải thích được kia chỉ đang phát tín hiệu vào chỗ
trống~\autocite{p24humangap}.

Lập luận ấy mạnh hơn một phép loại suy, vì nó đúng theo định nghĩa của cái
đang được tuyên bố. Định nghĩa~\ref{def:ch01-kha-nang-dien-giai} đặt khả năng
diễn giải là việc trình bày sao cho con người hiểu được, tức là một phát biểu
có hai đầu, và một phát biểu hai đầu không kiểm tra được nếu chỉ nhìn một đầu.
Điểm số của một chỉ số tự động là phát biểu về đầu bên phía mô hình. Nó có thể
đúng hoàn toàn mà đầu bên kia vẫn không xảy ra.

\index{nghiên cứu người dùng!nhãn do người gán}%
Có một dạng nghiên cứu nằm giữa hai mục ấy và bị chấm 0, và dạng này đáng tách
riêng vì nó dễ bị đọc nhầm thành có kiểm chứng. Đó là bài so kết quả của mình
với dữ liệu hoặc nhãn do người gán, chẳng hạn dùng nhãn người làm mốc, mà
không có bước đánh giá nào với người~\autocite{p24humangap}. Trong khung của
hình~\ref{fig:ch12-hai-nhanh-danh-gia}, một nghiên cứu như vậy vẫn nằm trọn ở
nhánh trên: người xuất hiện trong dữ liệu chứ không xuất hiện ở đầu ra của
nhánh dưới, nên phép đối chiếu nằm ngang giữa hai nhánh vẫn chưa được chạy.
Chương~\ref{ch:chi-so-khong-do-do-trung-thuc} đã gặp đúng cấu hình ấy ở một
chỗ khác: nhãn ground truth của bộ dữ liệu neo do thiết kế nhiệm vụ sinh ra và
được người kiểm lại, nhưng cái được chấm vẫn là chỉ số, không phải người.

Tình huống thứ nhất trong bốn tình huống bị chấm 0 mà
mục~\ref{sec:ch12-mot-cuoc-kiem-dem} liệt kê đáng quay lại ở đây, vì nó chạm
tới một ranh giới Phần~V sẽ dùng. Rudin phân biệt \enquote{explainable}, tức xử lý hậu
nghiệm một mô hình vốn là \tn{hộp đen}{black box}, với \enquote{interpretable},
tức tính chất sinh ra từ chính thiết kế mô hình dựa trên hiểu biết chuyên môn
về lĩnh vực. Một số bài bị chấm 0 đã dựng mô hình cùng với chuyên gia lĩnh
vực, nên thỏa định nghĩa thứ hai, nhưng lại không tiến hành đánh giá chính
thức nào với chính nhóm chuyên gia ấy hay với một nhóm rộng
hơn~\autocite{p24humangap}. Đây là trường hợp mà việc kiểm chứng đáng lẽ rẻ
nhất, vì người cần hỏi đã ngồi sẵn trong phòng, mà vẫn không được làm.

\index{nghiên cứu người dùng!lời kêu gọi của bài báo}%
Mục cuối bài là một lời kêu gọi, và nó dùng một phép so sánh với ngành dược.
Một công ty dược sẽ không bao giờ phân phối một loại thuốc mới mà bỏ qua thử
nghiệm lâm sàng diện rộng, kể cả khi họ thật sự tin vào hiệu quả của nó dựa
trên nguyên lý sinh học, trên bằng chứng mô phỏng và trên thí nghiệm trên
chuột; các tác giả hỏi vì sao mô hình XAI, thứ có thể tác động tới những quyết
định rủi ro cao, lại được phát hành mà chưa thử với người dùng của
nó~\autocite{p24humangap}.

Phép so sánh ấy có một phản bác quen thuộc, và bài báo trả lời nó ngay tại
chỗ. Phản bác là con người thì thiên lệch, và gom một mẫu người đại diện thì
khó. Câu trả lời của các tác giả là điều đó cũng đúng với mọi bộ dữ liệu chuẩn
mà học máy vẫn dùng, từ ImageNet tới tập bình luận phim IMDB tới các bộ hỏi
đáp y khoa: cả ba đều mang thiên lệch của con người từ cách chúng được sinh
ra, được tuyển chọn và từ bối cảnh văn hóa của chúng~\autocite{p24humangap}.
Lập luận ấy khép lại chỗ trốn cuối cùng: thiên lệch của con người không phải
lý do để không hỏi người, vì nó đã nằm sẵn trong những con số mà lĩnh vực này
vẫn báo cáo.

\index{nghiên cứu người dùng!chất lượng của phép kiểm chứng|)}%
